\documentclass[aspectratio=169, 12pt]{beamer}

% ── Thème et couleurs ────────────────────────────────────────────────
\usetheme{Madrid}
\usecolortheme{default}
\setbeamercolor{palette primary}{bg=suva_blue, fg=white}
\setbeamercolor{palette secondary}{bg=suva_mid, fg=white}
\setbeamercolor{palette tertiary}{bg=suva_dark, fg=white}
\setbeamercolor{palette quaternary}{bg=suva_dark, fg=white}
\setbeamercolor{structure}{fg=suva_blue}
\setbeamercolor{section in toc}{fg=suva_dark}
\setbeamercolor{block title}{bg=suva_blue, fg=white}
\setbeamercolor{block body}{bg=suva_blue!8}
\setbeamercolor{alerted text}{fg=danger_red}
\setbeamercolor{example text}{fg=done_green}

\setbeamertemplate{navigation symbols}{}
\setbeamertemplate{footline}{
  \leavevmode
  \hbox{%
    \begin{beamercolorbox}[wd=.33\paperwidth, ht=2.5ex, dp=1.125ex, left]{palette primary}%
      \hspace{1em}\insertshortauthor
    \end{beamercolorbox}%
    \begin{beamercolorbox}[wd=.34\paperwidth, ht=2.5ex, dp=1.125ex, center]{palette secondary}%
      \insertshorttitle
    \end{beamercolorbox}%
    \begin{beamercolorbox}[wd=.33\paperwidth, ht=2.5ex, dp=1.125ex, right]{palette primary}%
      \insertframenumber{} / \inserttotalframenumber\hspace{1em}
    \end{beamercolorbox}}%
}

% ── Packages ─────────────────────────────────────────────────────────
\usepackage[utf8]{inputenc}
\usepackage[T1]{fontenc}
\usepackage[french]{babel}
\usepackage{xcolor}
\usepackage{tcolorbox}
\tcbuselibrary{skins, breakable}
\usepackage{booktabs}
\usepackage{tikz}
\usetikzlibrary{arrows.meta, positioning, shapes.geometric}
\usepackage{fontawesome5}
\usepackage{hyperref}
\usepackage{pgfplots}
\pgfplotsset{compat=1.18}
\usepackage{colortbl}
\usepackage{array}

% ── Définitions couleurs ──────────────────────────────────────────────
\definecolor{suva_blue}{RGB}{30, 100, 180}
\definecolor{suva_mid}{RGB}{50, 130, 200}
\definecolor{suva_dark}{RGB}{20, 70, 130}
\definecolor{done_green}{RGB}{34, 139, 34}
\definecolor{done_bg}{RGB}{230, 255, 230}
\definecolor{danger_red}{RGB}{180, 30, 30}
\definecolor{warn_orange}{RGB}{200, 100, 0}
\definecolor{warn_bg}{RGB}{255, 248, 220}
\definecolor{gray_light}{RGB}{245, 245, 245}

% ── Notes speaker ─────────────────────────────────────────────────────
\setbeamertemplate{note page}{
  \vspace{0.5cm}
  {\color{suva_dark}\bfseries\large Ce que je dis :}\\[0.3cm]
  \insertnote
}

% ── Métadonnées ───────────────────────────────────────────────────────
\title[SuVa Protocol]{SuVa Protocol}
\subtitle{État d'avancement \& Prochaines étapes}
\author[M.L. MBENGUE]{Mohamed Lamine MBENGUE\\
  {\small Sécurité · DevOps · Dev Full Stack Web \& Mobile}\\
  {\footnotesize\texttt{mbenguemohamedlamine2022@gmail.com}}}
\date{Mars 2026}
\institute{\href{https://portfolio-alamine.web.app}{\texttt{portfolio-alamine.web.app}}}

% ═════════════════════════════════════════════════════════════════════
\begin{document}
% ═════════════════════════════════════════════════════════════════════

% ── SLIDE 1 : Titre ───────────────────────────────────────────────────
{
\setbeamercolor{background canvas}{bg=suva_dark}
\setbeamercolor{title}{fg=white}
\setbeamercolor{subtitle}{fg=suva_mid!60!white}
\setbeamercolor{author}{fg=white!90}
\setbeamercolor{date}{fg=white!70}
\setbeamercolor{institute}{fg=white!70}
\begin{frame}[plain]
  \vspace{1cm}
  \begin{center}
    {\color{white}\rule{0.6\textwidth}{1.5pt}}\\[0.5cm]
    {\color{white}\fontsize{28}{34}\selectfont\bfseries SuVa Protocol}\\[0.3cm]
    {\color{suva_mid!80!white}\large État d'avancement \& Prochaines étapes}\\[0.8cm]
    {\color{white}\rule{0.6\textwidth}{0.5pt}}\\[0.6cm]
    {\color{white!90}\normalsize Mohamed Lamine MBENGUE}\\
    {\color{white!70}\small Sécurité · DevOps · Dev Full Stack}\\[0.3cm]
    {\color{white!60}\footnotesize Mars 2026}
  \end{center}

  \note{
    \textit{Good morning / Good afternoon --- thanks for your time.}\\[0.2cm]
    Je m'appelle Mohamed Lamine MBENGUE, développeur spécialisé en sécurité et DevOps.
    Aujourd'hui je vais vous présenter l'état d'avancement du projet SuVa ---
    ce qui a été accompli, ce qui reste, et les questions clés pour la suite.\\[0.2cm]
    \textit{This should take about 10 to 15 minutes, then we'll have time for Q\&A.}
  }
\end{frame}
}

% ── SLIDE 2 : Agenda ──────────────────────────────────────────────────
\begin{frame}{Plan de la présentation}

  \begin{columns}[T]
    \begin{column}{0.5\textwidth}
      \begin{tcolorbox}[colback=suva_blue!6, colframe=suva_blue!40,
        title={\bfseries Au programme}, fonttitle=\color{suva_dark}]
        \begin{enumerate}
          \item[\faCheckCircle] Qu'est-ce que SuVa ?
          \item[\faCheckCircle] Ce qui est \textbf{déjà fait}
          \item[\faArrowCircleRight] Ce qui \textbf{reste à faire}
          \item[\faCalendar] Feuille de route
          \item[\faQuestion] Questions \& échanges
        \end{enumerate}
      \end{tcolorbox}
    \end{column}
    \begin{column}{0.45\textwidth}
      \begin{tcolorbox}[colback=done_bg, colframe=done_green,
        title={\bfseries\color{done_green} Avancement global}]
        \begin{center}
          \begin{tikzpicture}
            \fill[done_bg, draw=done_green, rounded corners] (0,0) rectangle (4*5/12, 0.5);
            \fill[gray!20, draw=gray!40, rounded corners] (4*5/12, 0) rectangle (4, 0.5);
            \node[font=\small\bfseries] at (2, 0.25) {5 / 12 composants};
          \end{tikzpicture}\\[0.2cm]
          {\color{done_green}\bfseries\Large 41\%} terminé
        \end{center}
      \end{tcolorbox}
    \end{column}
  \end{columns}

  \note{
    \textit{Quick overview of what we'll cover.}\\[0.2cm]
    On commence par rappeler ce qu'est SuVa --- \textit{making sure we're all on the same page}.
    Ensuite je vous montre ce qui est fonctionnel aujourd'hui.
    Puis les prochaines étapes prioritaires.
    Et on termine sur les questions que j'ai pour vous --- \textit{your answers will shape the next phase.}
  }
\end{frame}

% ── SLIDE 3 : Vision ──────────────────────────────────────────────────
\begin{frame}[shrink=8]{Qu'est-ce que SuVa ?}

  \begin{tcolorbox}[colback=suva_blue!8, colframe=suva_blue,
    title={\bfseries SuVa (Super Value) Protocol}, fonttitle=\color{white},
    top=3pt, bottom=3pt]
    \small Un \textbf{stablecoin quantique-résistant} pour blockchains Ethereum.\\
    Les ordinateurs quantiques menacent ECDSA (la signature actuelle d'Ethereum).
    SuVa remplace ECDSA par des signatures \textbf{post-quantiques standardisées par le NIST}.
  \end{tcolorbox}

  \vspace{0.2cm}

  \begin{columns}[T]
    \begin{column}{0.31\textwidth}
      \begin{tcolorbox}[colback=gray_light, colframe=suva_blue!40, height=3cm]
        \centering
        {\color{suva_blue}\Large\faShieldAlt}\\[0.15cm]
        \textbf{Sécurité PQ}\\[0.05cm]
        {\small Dilithium \& Falcon\\(NIST FIPS 204/206)}
      \end{tcolorbox}
    \end{column}
    \begin{column}{0.31\textwidth}
      \begin{tcolorbox}[colback=gray_light, colframe=suva_blue!40, height=3cm]
        \centering
        {\color{suva_blue}\Large\faExchangeAlt}\\[0.15cm]
        \textbf{Wrapping 1:1}\\[0.05cm]
        {\small USDT $\to$ sUSDT\\USDC $\to$ sUSDC}
      \end{tcolorbox}
    \end{column}
    \begin{column}{0.31\textwidth}
      \begin{tcolorbox}[colback=gray_light, colframe=suva_blue!40, height=3cm]
        \centering
        {\color{suva_blue}\Large\faEthereum}\\[0.15cm]
        \textbf{EVM Ready}\\[0.05cm]
        {\small Ethereum, L2\\ERC-4337 ready}
      \end{tcolorbox}
    \end{column}
  \end{columns}

  \note{
    \textit{So, the big picture first.}\\[0.2cm]
    SuVa, c'est simple à expliquer en une phrase : c'est un stablecoin qui survivra à l'ère quantique.\\[0.2cm]
    Aujourd'hui, quand vous envoyez des fonds sur Ethereum, vous signez avec ECDSA.
    Un ordinateur quantique suffisamment puissant pourra casser ECDSA ---
    \textit{experts estimate within 10 to 20 years}.\\[0.2cm]
    SuVa remplace ECDSA par Dilithium ou Falcon, les nouveaux standards officiels du NIST 2024.\\[0.2cm]
    \textit{Bottom line :} vous déposez vos USDT dans SuVa, vous recevez des sUSDT sécurisés.
    Chaque retrait ou transfert nécessite une signature post-quantique.
  }
\end{frame}

% ── SLIDE 4 : Architecture ────────────────────────────────────────────
\begin{frame}{Architecture du système}

  \begin{center}
  \begin{tikzpicture}[
    node distance=1.6cm,
    box/.style={draw, rounded corners, minimum width=2.8cm, minimum height=1cm,
                align=center, font=\small\bfseries},
    arrow/.style={-Stealth, thick, suva_blue},
    label_style/.style={font=\tiny, color=gray},
  ]
    \node[box, fill=done_bg, draw=done_green] (user) {\faUser\\Utilisateur};
    \node[box, fill=suva_blue!15, draw=suva_blue, right=1.8cm of user] (python)
      {\faPython\\Python off-chain\\{\tiny keygen + sign}};
    \node[box, fill=suva_blue!15, draw=suva_blue, right=1.8cm of python] (sol)
      {\faFileCode\\Solidity on-chain\\{\tiny verify + vault}};
    \node[box, fill=gray!15, draw=gray!60, above=1.2cm of python] (sdk)
      {\faCode\\SDK\\{\tiny (à venir)}};
    \node[box, fill=warn_bg, draw=warn_orange, below=1.2cm of sol] (vault)
      {\faLock\\Vault Contract\\{\tiny sUSDT/sUSDC}};

    \draw[arrow] (user) -- (python) node[midway, above, label_style] {demande};
    \draw[arrow] (python) -- (sol) node[midway, above, label_style] {sig hex};
    \draw[arrow] (sol) -- (vault) node[midway, right, label_style] {si valide};
    \draw[arrow, dashed, gray] (sdk) -- (python);
    \draw[arrow, dashed, gray] (sdk) -- (sol);

    % Légende
    \node[font=\tiny, done_green] at (0, -1.8) {\faCheckCircle\ Fait};
    \node[font=\tiny, warn_orange] at (1.5, -1.8) {\faTools\ À faire};
    \node[font=\tiny, gray] at (3, -1.8) {- - - Futur};
  \end{tikzpicture}
  \end{center}

  \note{
    \textit{Here's how the system works, end to end.}\\[0.2cm]
    L'utilisateur utilise le SDK pour générer ses clés et signer --- \textit{all off-chain, on his device, free.}\\[0.2cm]
    La signature est envoyée au contrat Solidity sur Ethereum qui vérifie.
    \textit{If valid}, le vault libère les fonds.\\[0.2cm]
    En vert : déjà fait. En orange : \textit{work in progress.}
  }
\end{frame}

% ── SLIDE 5 : Ce qui est fait ─────────────────────────────────────────
\begin{frame}[shrink=10]{Ce qui est \textcolor{done_green}{déjà fait} \faCheckCircle}

  \begin{columns}[T]
    \begin{column}{0.48\textwidth}
      \begin{block}{\faCode\ Couche Python (off-chain)}
        \small
        \begin{itemize}\setlength\itemsep{1pt}
          \item[\checkmark] Keygen, Sign, Verify
          \item[\checkmark] Dilithium2 / 3 / 5 (NIST)
          \item[\checkmark] ETHDilithium (variante Keccak)
          \item[\checkmark] Tests KAT officiels NIST \textcolor{done_green}{\bfseries 4/4}
        \end{itemize}
      \end{block}
      \vspace{0.15cm}
      \begin{block}{\faEthereum\ Solidity (on-chain)}
        \small
        \begin{itemize}\setlength\itemsep{1pt}
          \item[\checkmark] Vérificateur Dilithium2
          \item[\checkmark] NTT optimisée (Yul)
          \item[\checkmark] PRNG Keccak compatible
          \item[\checkmark] Tests croisés Python$\leftrightarrow$Solidity
        \end{itemize}
      \end{block}
    \end{column}
    \begin{column}{0.48\textwidth}
      \begin{tcolorbox}[colback=done_bg, colframe=done_green,
        title={\bfseries\color{done_green} Chiffres clés},
        top=3pt, bottom=3pt]
        \renewcommand{\arraystretch}{1.2}
        \small
        \begin{tabular}{lr}
          \textbf{Tests qui passent} & \textcolor{done_green}{\bfseries 4/4} \\
          \textbf{Gas Dilithium2} & $\sim$13.5M \\
          \textbf{Gas ETHDilithium} & $\sim$6.6M \\
          \textbf{Sig. Dilithium2} & 2\,420 oct. \\
          \textbf{Standard} & NIST FIPS 204 \\
        \end{tabular}
      \end{tcolorbox}
      \vspace{0.15cm}
      \begin{tcolorbox}[colback=warn_bg, colframe=warn_orange,
        title={\bfseries Problème : le gas},
        top=3pt, bottom=3pt]
        {\small 6.6M gas = \textbf{100x trop cher} vs ECDSA (3\,000 gas).\\
        La solution : \textbf{Falcon}.}
      \end{tcolorbox}
    \end{column}
  \end{columns}

  \note{
    \textit{Let me walk you through what's already done.}\\[0.2cm]
    Du côté Python : Dilithium tourne correctement, keygen, sign, verify.
    On a validé avec les vecteurs officiels du NIST --- \textit{all 4 tests pass}.\\[0.2cm]
    Du côté Solidity : le vérificateur est opérationnel sur l'EVM.
    Il y a une optimisation en assembleur Yul pour réduire le gas.\\[0.2cm]
    \textit{But here's the catch} --- 6.6 millions de gas c'est 100 fois trop cher pour un usage réel.
    ECDSA coûte 3000 gas. Falcon coûterait 350 000 gas.
    C'est pour ça que Falcon est la \textit{top priority}.
  }
\end{frame}

% ── SLIDE 6 : Comparaison Dilithium vs Falcon ─────────────────────────
\begin{frame}[shrink=8]{Pourquoi passer à Falcon ? \faArrowRight\ 40x moins cher}

  \begin{center}
  \renewcommand{\arraystretch}{1.25}
  \small
  \begin{tabular}{lccccc}
    \toprule
    \textbf{Schéma} & \textbf{Gas verify} & \textbf{vs ECDSA} & \textbf{Signature} & \textbf{Standard} & \textbf{Statut} \\
    \midrule
    ECDSA (actuel) & 3\,000 & $\times$1 & 65 oct. & --- & \textcolor{danger_red}{Vulnérable PQ} \\
    Dilithium2 & $\sim$13.5M & $\times$4\,500 & 2\,420 oct. & FIPS 204 & \textcolor{done_green}{\checkmark Fait} \\
    ETHDilithium & $\sim$6.6M & $\times$2\,200 & 2\,420 oct. & --- & \textcolor{done_green}{\checkmark Fait} \\
    \rowcolor{done_bg}
    \textbf{Falcon-512} & $\sim$\textbf{350K} & $\times$\textbf{117} & \textbf{666 oct.} & \textbf{FIPS 206} & \textcolor{warn_orange}{\bfseries Priorité \#1} \\
    SPHINCS+ & 50--500M & $\times$166K+ & 8--50 Ko & FIPS 205 & \textcolor{gray}{Plus tard} \\
    \bottomrule
  \end{tabular}
  \end{center}

  \vspace{0.15cm}
  \begin{columns}
    \begin{column}{0.48\textwidth}
      \begin{tcolorbox}[colback=done_bg, colframe=done_green]
        \centering
        {\large\bfseries\color{done_green} Falcon-512}\\
        350K gas · 666 octets\\
        {\small Réduction de \textbf{95\%} vs ETHDilithium}
      \end{tcolorbox}
    \end{column}
    \begin{column}{0.48\textwidth}
      \begin{tcolorbox}[colback=suva_blue!8, colframe=suva_blue]
        \centering
        ZKNox a déjà une implémentation\\
        \textbf{ETHFALCON} disponible\\
        {\small github.com/ZKNox/ETHFALCON}
      \end{tcolorbox}
    \end{column}
  \end{columns}

  \note{
    \textit{This table tells the whole story.}\\[0.2cm]
    ECDSA : 3000 gas, la référence. Mais vulnérable aux ordinateurs quantiques.\\[0.2cm]
    Dilithium : ça marche, mais 6.6 millions de gas c'est vraiment trop cher
    pour chaque transaction. \textit{Not production-ready.}\\[0.2cm]
    Falcon, c'est la cible : 350 000 gas, \textit{that's 40x cheaper} que ETHDilithium.
    La signature est aussi plus petite : 666 octets contre 2420.\\[0.2cm]
    \textit{Good news} : ZKNox a déjà une implémentation Falcon pour l'EVM --- ETHFALCON.
    \textit{We don't start from scratch.} Notre travail : intégrer ça dans SuVa.
  }
\end{frame}

% ── SLIDE 7 : Prochaines étapes ───────────────────────────────────────
\begin{frame}[shrink=15]{Ce qui reste à faire --- Priorités}

  \begin{tcolorbox}[colback=danger_red!8, colframe=danger_red,
    title={\bfseries\color{white} \faExclamationTriangle\ Priorité \#1 : Intégration Falcon (ETHFALCON)},
    fonttitle=\color{white}, top=3pt, bottom=3pt]
    \small
    \begin{enumerate}\setlength\itemsep{1pt}
      \item Étudier et adapter les contrats ETHFALCON de ZKNox
      \item Implémenter le signing Falcon en Python (off-chain)
      \item Tests croisés Python $\to$ Solidity
      \item Intégrer dans le vault SuVa
    \end{enumerate}
  \end{tcolorbox}

  \vspace{0.15cm}

  \begin{columns}[T]
    \begin{column}{0.31\textwidth}
      \begin{tcolorbox}[colback=warn_bg, colframe=warn_orange,
        title={\bfseries\small Priorité \#2}, top=2pt, bottom=2pt]
        {\small \faLock\ \textbf{Vault Contract}\\
        Wrap USDT $\to$ sUSDT\\
        Deposit / Withdraw}
      \end{tcolorbox}
    \end{column}
    \begin{column}{0.31\textwidth}
      \begin{tcolorbox}[colback=warn_bg, colframe=warn_orange,
        title={\bfseries\small Priorité \#2}, top=2pt, bottom=2pt]
        {\small \faWallet\ \textbf{ERC-4337 Wallet}\\
        Smart wallet PQ\\
        Replaces ECDSA}
      \end{tcolorbox}
    \end{column}
    \begin{column}{0.31\textwidth}
      \begin{tcolorbox}[colback=gray!10, colframe=gray!50,
        title={\bfseries\small Priorité \#3}, top=2pt, bottom=2pt]
        {\small \faCode\ \textbf{SDK + Audit}\\
        Python \& JS SDK\\
        Security audit}
      \end{tcolorbox}
    \end{column}
  \end{columns}

  \note{
    \textit{Here's the roadmap for what's next.}\\[0.2cm]
    La priorité absolue c'est Falcon. C'est le \textit{blocker} pour tout le reste.
    Sans un coût gas acceptable, pas de déploiement en production.\\[0.2cm]
    Ensuite le Vault Contract --- le cœur financier du produit, \textit{where the money lives}.\\[0.2cm]
    Le wallet ERC-4337 permet aux utilisateurs d'utiliser une signature post-quantique
    sans modifier le protocole Ethereum. \textit{That's the beauty of account abstraction.}\\[0.2cm]
    Et enfin le SDK et l'audit de sécurité avant le \textit{mainnet launch}.
  }
\end{frame}

% ── SLIDE 8 : Roadmap ─────────────────────────────────────────────────
\begin{frame}[shrink=8]{Feuille de route --- Planning estimé}

  \begin{center}
  \begin{tikzpicture}[x=2.2cm, y=1.1cm]
    % Axe
    \draw[-Stealth, thick, suva_blue] (-0.2, 0) -- (5.2, 0);

    % Phase 1 : FAIT
    \fill[done_bg, draw=done_green, rounded corners] (0, 0.2) rectangle (1.6, 0.9);
    \node[font=\small\bfseries, color=done_green] at (0.8, 0.62) {Phase 1 \checkmark};
    \node[font=\tiny, color=done_green] at (0.8, 0.33) {Dilithium + NTT};

    % Phase 2 : EN COURS / PRIORITÉ
    \fill[warn_bg, draw=warn_orange, rounded corners] (1.6, 0.2) rectangle (3, 0.9);
    \node[font=\small\bfseries, color=warn_orange] at (2.3, 0.62) {Phase 2};
    \node[font=\tiny, color=warn_orange] at (2.3, 0.33) {Falcon (8--10 sem.)};

    % Phase 3
    \fill[gray!10, draw=gray!50, rounded corners] (3, 0.2) rectangle (4.1, 0.9);
    \node[font=\small\bfseries, color=gray!70] at (3.55, 0.62) {Phase 3};
    \node[font=\tiny, color=gray] at (3.55, 0.33) {Vault + ERC-4337};

    % Phase 4
    \fill[gray!10, draw=gray!50, rounded corners] (4.1, 0.2) rectangle (5, 0.9);
    \node[font=\small\bfseries, color=gray!70] at (4.55, 0.62) {Phase 4};
    \node[font=\tiny, color=gray] at (4.55, 0.33) {SDK + Audit};

    % Marqueurs temps
    \foreach \x/\label in {0/Déc 2024, 1.6/Mars 2026, 3/Juin 2026, 4.1/Sept 2026, 5/Nov 2026} {
      \draw (\x, -0.1) -- (\x, 0.1);
      \node[below, font=\tiny, rotate=15] at (\x, -0.12) {\label};
    }

    % Indicateur "Vous êtes ici"
    \draw[danger_red, very thick, -Stealth] (1.6, 1.3) -- (1.6, 0.95);
    \node[font=\tiny\bfseries, color=danger_red] at (1.6, 1.42) {Nous sommes ici};
  \end{tikzpicture}
  \end{center}

  \vspace{0.05cm}
  \small\renewcommand{\arraystretch}{1.2}
  \begin{tabular}{clp{7cm}c}
    \toprule
    \textbf{Phase} & \textbf{Contenu} & \textbf{Key deliverables} & \textbf{Durée} \\
    \midrule
    \rowcolor{done_bg}
    1 & Fondations & NTT, Dilithium Python + Solidity, Tests KAT & Fait \\
    \rowcolor{warn_bg}
    2 & \textbf{Falcon} & ETHFALCON adapté, Python signing, Tests croisés & 8--10 sem. \\
    3 & Vault + 4337 & Vault Contract, ERC-4337 Smart Wallet & 8--12 sem. \\
    4 & SDK + Audit & SDK Python/JS, Security audit, Mainnet deploy & 8 sem. \\
    \bottomrule
  \end{tabular}

  \note{
    \textit{Here's the timeline I consider realistic.}\\[0.2cm]
    La phase 1 est terminée. \textit{We are right here} --- début de la phase Falcon.\\[0.2cm]
    \textit{If we start now}, Falcon opérationnel vers juin 2026.
    Ensuite 3 mois pour le vault et le wallet ERC-4337.
    Et fin 2026, \textit{mainnet deployment} après audit.\\[0.2cm]
    Ce sont des estimations --- elles dépendent de vos réponses et des ressources disponibles.
    \textit{Flexible, but realistic.}
  }
\end{frame}

% ── SLIDE 9 : Questions ───────────────────────────────────────────────
\begin{frame}[shrink=8]{Questions pour avancer}

  \begin{columns}[T]
    \begin{column}{0.48\textwidth}
      \begin{block}{\faCog\ Architecture}
        \small\begin{itemize}\setlength\itemsep{1pt}
          \item Falcon seul, ou Dilithium + Falcon en parallèle ?
          \item Vault upgradeable (proxy) ou immutable ?
          \item Cible réseau : Ethereum mainnet ou L2 d'abord ?
        \end{itemize}
      \end{block}
      \vspace{0.1cm}
      \begin{block}{\faShieldAlt\ Sécurité}
        \small\begin{itemize}\setlength\itemsep{1pt}
          \item ETHFALCON utilisé tel quel ou modifié ?
          \item Mécanisme de récupération si clé PQ perdue ?
          \item Audit prévu, et par qui ?
        \end{itemize}
      \end{block}
    \end{column}
    \begin{column}{0.48\textwidth}
      \begin{block}{\faGasPump\ Gas \& Budget}
        \small\begin{itemize}\setlength\itemsep{1pt}
          \item Budget gas max acceptable par transaction ?
          \item Date cible testnet ? Mainnet ?
        \end{itemize}
      \end{block}
      \vspace{0.1cm}
      \begin{block}{\faUsers\ Usage}
        \small\begin{itemize}\setlength\itemsep{1pt}
          \item Scénario : wallet mobile, hardware ou serveur ?
          \item Comment l'utilisateur stocke-t-il sa clé PQ ?
        \end{itemize}
      \end{block}
      \vspace{0.1cm}
      \begin{tcolorbox}[colback=suva_blue!8, colframe=suva_blue,
        top=3pt, bottom=3pt]
        {\small Ces réponses permettent de \textbf{prioriser} et d'éviter de retravailler
        les mêmes composants deux fois.}
      \end{tcolorbox}
    \end{column}
  \end{columns}

  \note{
    \textit{These are the key questions I need answered to move forward efficiently.}\\[0.2cm]
    La plus importante : Falcon seul en production, ou Dilithium + Falcon en parallèle ?
    \textit{This impacts the whole architecture.}\\[0.2cm]
    Deuxième : le réseau cible. Mainnet directement, ou d'abord un L2 ?
    Sur L2, le gas est 100x moins cher --- \textit{Dilithium might already be viable there.}\\[0.2cm]
    Et le budget gas : quel est le seuil acceptable ?
    Moins de 500K gas ? \textit{That guides our optimization efforts directly.}\\[0.2cm]
    \textit{I have a full document with all these questions if you want to review after the call.}
  }
\end{frame}

% ── SLIDE 10 : Conclusion ─────────────────────────────────────────────
{
\setbeamercolor{background canvas}{bg=suva_dark}
\begin{frame}[plain]
  \begin{center}
    \vspace{0.8cm}
    {\color{white}\rule{0.5\textwidth}{1pt}}\\[0.5cm]

    \begin{tcolorbox}[colback=white!10!suva_dark, colframe=white!30,
      width=0.85\textwidth]
      \centering
      {\color{white}\bfseries\large Résumé en 3 points}\\[0.3cm]
      \begin{tabular}{cl}
        \textcolor{done_green}{\Large\checkmark} & {\color{white} Dilithium est fonctionnel et validé} \\[0.2cm]
        \textcolor{warn_orange}{\Large\faArrowRight} & {\color{white} Falcon est la priorité absolue (350K gas)} \\[0.2cm]
        \textcolor{suva_mid}{\Large\faQuestion} & {\color{white} Vos réponses guident la prochaine phase} \\
      \end{tabular}
    \end{tcolorbox}

    % \vspace{0.5cm}
    {\color{white}\rule{0.5\textwidth}{0.5pt}}\\[0.5cm]

    {\color{white!90}\bfseries Mohamed Lamine MBENGUE}\\[0.1cm]
    {\color{white!70}\small +221 78 178 44 36}\\
    {\color{white!70}\small mbenguemohamedlamine2022@gmail.com}\\[0.1cm]
    {\color{suva_mid}\small\href{https://portfolio-alamine.web.app}{\texttt{portfolio-alamine.web.app}}}

    \vspace{0.3cm}
    {\color{white!60}\large Merci --- Des questions ?}
  \end{center}

  \note{
    \textit{To wrap it up --- three key points.}\\[0.2cm]
    Premièrement : la fondation est solide. Dilithium est implémenté, testé, validé.
    \textit{The code works on both sides, Python and Solidity.}\\[0.2cm]
    Deuxièmement : la priorité c'est Falcon. \textit{It's the only scheme with acceptable gas cost for production.}
    ZKNox a déjà la base, \textit{we can build on top of that.}\\[0.2cm]
    Troisièmement : vos réponses vont directement influencer la phase suivante.
    \textit{That will save us time and money.}\\[0.2cm]
    \textit{Thank you for your time --- I'm happy to take any questions.}
  }
\end{frame}
}

% ── SLIDE BONUS : Notes speaker (à ne pas projeter) ───────────────────
\appendix

\begin{frame}[noframenumbering]{Guide speaker --- Ce que je dis slide par slide}
  \begin{tcolorbox}[colback=gray!5, colframe=gray!40, breakable,
    title={\bfseries Instructions pour utiliser les notes}]
    \begin{itemize}
      \item Ouvrir le PDF en mode \textbf{Présentation Dual Screen} (ex. LibreOffice Impress,
            Adobe Acrobat, ou \texttt{pdfpc})
      \item Les notes apparaissent sur ton écran, pas sur le projector
      \item Avec \textbf{pdfpc} (Linux/Mac) : \texttt{pdfpc 10\_presentation\_client.pdf}
      \item Avec \textbf{Google Slides} : importer le PDF, ajouter les notes manuellement
            depuis ce document
      \item \textbf{Pour une présentation en ligne (Zoom/Meet)} : activer le partage d'écran
            uniquement sur la fenêtre du PDF, et garder les notes sur ton autre écran
    \end{itemize}
  \end{tcolorbox}

  % \vspace{0.3cm}
  % {\bfseries Durée estimée par slide :}
  % \begin{center}
  % \begin{tabular}{clc}
  %   \toprule
  %   \textbf{Slide} & \textbf{Contenu} & \textbf{Durée} \\
  %   \midrule
  %   1 & Titre & 30 sec \\
  %   2 & Agenda & 45 sec \\
  %   3 & Vision SuVa & 2 min \\
  %   4 & Architecture & 1 min 30 \\
  %   5 & Ce qui est fait & 2 min \\
  %   6 & Falcon vs Dilithium & 2 min \\
  %   7 & Prochaines étapes & 1 min 30 \\
  %   8 & Roadmap & 1 min 30 \\
  %   9 & Questions & 2 min + échange \\
  %   10 & Conclusion & 1 min \\
  %   \midrule
  %   & \textbf{Total} & \textbf{$\sim$15 min} \\
  %   \bottomrule
  % \end{tabular}
  % \end{center}
\end{frame}

\end{document}
